% Generated draft: requires mathematical and source review.
\documentclass[11pt]{article}
\usepackage[margin=25mm]{geometry}
\usepackage{fontspec}
\setmainfont{Latin Modern Roman}
\usepackage{amsmath,amssymb,mathtools,unicode-math}
\setmathfont{Latin Modern Math}
\usepackage{enumitem,needspace,xurl,hyperref}
\setlength{\parindent}{0pt}
\setlength{\parskip}{5pt}
\setlength{\emergencystretch}{3em}
\newcommand{\N}{\mathbb{N}}
\newcommand{\Z}{\mathbb{Z}}
\newcommand{\Q}{\mathbb{Q}}
\newcommand{\R}{\mathbb{R}}
\newcommand{\C}{\mathbb{C}}
\newcommand{\F}{\mathbb{F}}
\newcommand{\A}{\mathbb{A}}
\newcommand{\PP}{\mathbb{P}}
\newcommand{\E}{\mathbb{E}}
\newcommand{\eps}{\varepsilon}
\newcommand{\dd}{\,\mathrm{d}}
\begin{document}
\section*{UnsolvedMath ID: 30004643}
\textbf{Source problem code:} catalogue entry
\subsection{Definitions and mathematical statement}
Let \(T\) be transcendental over \(\Q\), and let \(E/\Q(T)\) be a finite Galois extension with Galois group \(G\). The extension is \(\Q\)-regular if \(E\cap\overline{\Q}=\Q\). Let \(C/\Q\) be the smooth projective geometrically integral curve with function field \(E\), and write \(g=g(C)\). Assume \(g\geq 2\).

The recorded question asks whether, for every such extension,
\[
\lim_{x\to\infty}
\frac{\bigl|\operatorname{Sp}(E)\cap S(G,x)\bigr|}
     {\bigl|S(G,x)\bigr|}=0.
\]
Here the record calls \(\operatorname{Sp}(E)\) the specialization set and describes the limit as density zero in \(S(G)\), but does not define these sets or the ordering \(S(G,x)\).

A standard interpretation of \(S(G)\) is the set of \(\Q\)-isomorphism classes of finite Galois extensions of \(\Q\) with group isomorphic to \(G\). Under that interpretation, \(\operatorname{Sp}(E)\cap S(G)\) consists of classes obtained from residue extensions at unramified rational specialization points with full Galois group \(G\). These conventions require confirmation from the source; in particular, the supplied record does not specify the treatment of the point at infinity or of specializations with smaller Galois group.

Most importantly, no height, discriminant bound, or other rule defining \(S(G,x)\) is supplied. The displayed limit therefore retains the source notation without selecting an ordering. A precise density assertion requires that rule and a counting convention, with \(S(G,x)\) finite and its cardinality nonzero for all sufficiently large \(x\). These are unresolved formulation requirements, not additional hypotheses asserted by the source.
\subsection{Short English statement}
For every \(\Q\)-regular Galois cover of function-field genus at least two, must its specializations have density zero among extensions of \(\Q\) with the same Galois group? The supplied record leaves the counting order and some specialization conventions undefined, so the density statement requires source clarification.
\subsection{Sources}
\begin{itemize}
\item Local UnsolvedMath registry: \texttt{data/unsolvedmath/problems.json}, permanent ID \(30004643\), title ``Density of Specializations of High-Genus Galois Covers,'' source problem code \texttt{catalogue entry}. The record gives proposed year 2021 and supplies no individual proposer.
\item Supplied report citation: \emph{Homotopic and Geometric Galois Theory (2022)}, Oberwolfach Reports, pp.\ 663--744, DOI \texttt{10.4171/owr/2021/12}. \url{https://doi.org/10.4171/owr/2021/12}. The literature list also labels this source ``Oberwolfach Report 12/2021.'' Both supplied year designations are retained without resolving their relationship.
\item Supplied literature reference: ``Density results for specialization sets.'' \url{https://doi.org/10.1017/S1474748019000537}.
\item Supplied literature assessment, checked on exactly \texttt{2026-08-22}: the status is labeled ``partially solved.'' The assessment says that density zero is proved conditionally for sufficiently many branch points and in several unconditional families, while the genus-at-least-two question for a general cover remains unresolved. Its listed sources are \url{https://doi.org/10.1017/S1474748019000537} and \url{https://doi.org/10.4171/owr/2021/12}. This is supplied dated curation derived from the OWR material, not an independently verified literature assessment; its time-sensitive status and cited claims require independent verification.
\end{itemize}
\end{document}
